% front/series-preface.tex —— 原书丛书序（Preface to the Series）
% 作者：Shi Zhongci，2005 年 7 月。
% 注：每段必须写成一行（不得硬换行），否则 xeCJK 会把换行当成空格。
\cleardoublepage
\phantomsection
\chapter*{《信息与计算科学丛书》序}
\addcontentsline{toc}{chapter}{《信息与计算科学丛书》序}
\markboth{《信息与计算科学丛书》序}{《信息与计算科学丛书》序}

20 世纪 70 年代以来, 科学出版社已陆续出版三十余种「计算方法丛书」. 这套丛书由已故中国科学院院士, 中国科学院计算中心主任冯康先生倡导并主持. 丛书及时反映了计算数学的前沿方向与最新研究成果, 对青年科学工作者和整个研究群体产生了很大影响, 对我国计算数学的发展起了重要作用.

为适应科学发展的新形势, 1998 年, 中华人民共和国教育部将计算数学, 数值算法, 信息科学, 运筹学与控制论等学科合并为「信息与计算科学」这一新学科. 科学出版社亦相应调整了丛书编委会, 并将丛书更名为「信息与计算科学丛书」. 第一届编委会会议于 2004 年 9 月在北京召开, 讨论了新丛书的目标, 方向与内容.

新丛书旨在向高年级本科生, 研究生和从事这些领域研究的科学工作者介绍信息与计算科学的现状与前沿, 因而将以具体而系统的方式阐述信息与计算科学的进展, 并兼及相关交叉学科的发展.

谨向历任编委会成员与编辑部同仁, 以及所有为「计算方法丛书」作出重要贡献的数学家致谢. 正是由于他们的努力, 这套丛书在学术界赢得了卓越声誉. 我真诚希望各位继续支持「信息与计算科学丛书」, 使新丛书同样能推动本世纪信息与计算科学的发展.

\vspace{1.5em}
\rightline{石钟慈}
\rightline{2005 年 7 月}
